% =============================================================================
%   I. INTRODUCTION
% =============================================================================
\section{Introduction}
\label{sec:intro}

ESG investing now sits at institutional scale: global ESG assets under
management exceed \$35.3~trillion, roughly one-third of professionally managed
capital \cite{gsia_2022}. At the same time, the core signal used in many
allocation pipelines remains unstable across providers. Berg et al.
document substantial ESG rating divergence, with average cross-provider
correlation around $r=0.54$ \cite{berg_2022}. This tension is most acute in
mid-cap equities, where coverage is thinner than in large caps but portfolio
capacity and investability remain materially higher than in small caps. The
segment is economically important, yet methodologically underserved.

This paper addresses a concrete gap: there is no transparent and reproducible
ESG-integrated nine-factor index for mid-cap equities built on publicly
available data and explicit design choices. Existing products and studies often
mix proprietary inputs, opaque weighting rules, and limited methodological
disclosure, which makes independent verification difficult and weakens
comparability across contexts.

We focus on mid-caps because they are the practical ``sweet spot'' for ESG
integration tests: large enough for institutional portfolio capacity, yet small
enough that ESG disclosure and coverage gaps still materially affect investment
decisions. In aggregate, mid-caps represent roughly one-fifth of global equity
market capitalization but receive relatively few studies in ESG research.
Economically, they offer a
distinct trade-off---typically higher growth optionality than large-caps and
better investability than small-caps. Methodologically, this is precisely where
the data gap is most binding: large-caps increasingly face mandatory ESG
reporting, while many small-caps are excluded from investable universes;
mid-caps fall between these extremes. This segment therefore provides the most
actionable setting to test whether transparent ESG integration can work without
proprietary vendor datasets.

Prior work establishes why this gap matters, but also why standard approaches
are insufficient. First, provider disagreement creates measurement noise that
is not a minor implementation detail but a structural problem for ESG-based
allocation \cite{berg_2022}. Second, broad evidence links ESG and financial
outcomes on average \cite{friede_2015}, while materiality-based evidence shows
that not all ESG indicators are equally decision-relevant across sectors
\cite{khan_2016}. Third, factor-investing research shows that expected-return
variation is multi-dimensional rather than single-metric \cite{fama_french_2015},
so ESG-only or financial-only screens are likely to discard useful information.
Fourth, index construction is inherently sensitive to weighting and aggregation
choices \cite{oecd_2008}; black-box specifications therefore limit scientific
reproducibility. Finally, recent equilibrium frameworks suggest that ESG and
climate preferences can become priced in expected returns and cost of capital
\cite{pastor_2021, pedersen_2021}, making transparent integration design
especially relevant for both research and practice.

We therefore ask three research questions.
\begin{itemize}
\item[\textbf{RQ1.}] Can a transparent nine-factor index combining ESG and
financial metrics identify mid-cap equities with superior risk-adjusted
cross-sectional performance proxies?
\item[\textbf{RQ2.}] Does the integration of ESG factors provide incremental
value beyond traditional financial metrics?
\item[\textbf{RQ3.}] How robust is the index to methodological choices
(weighting, composition, market regime)?
\end{itemize}

Our approach is a nine-factor composite model tailored to a 276-company mid-cap
universe (186 US, 90 India) across 11 GICS sectors using 202 variables. The
final framework combines ESG, financial, risk-adjusted, growth,
sector-position, stability, operational, market, and value dimensions under
explicit and reproducible weighting rules. We remove
\texttt{similarity\_rank} from the deployed specification because it is
collinear with ESG ($r=0.70$). Within the ESG block, we apply SASB-informed
sector materiality weighting with pillar weights $E=0.20$, $S=0.35$, and
$G=0.45$ rather than uniform treatment of sector-specific exposures. The
deployed balanced profile assigns weights of 22\% ESG, 18\% financial, 15\%
risk-adjusted, 14\% growth, 8\% sector position, 8\% stability, 6\%
operational, 5\% market, and 4\% value. We implement three investor profiles
(ESG-first, balanced, and financial-first) to make preference trade-offs
explicit rather than implicit. Because mid-cap ESG data coverage is incomplete
in open datasets, we use a hybrid data strategy that combines real Yahoo
Finance and SEC EDGAR governance inputs with proxy-derived environmental and
social indicators plus provenance-level tracking. We also include circularity
correction in validation design so that evaluation does not mechanically reward
overlapping signals.

This paper makes four contributions:
\begin{itemize}
\item We introduce a \textbf{reproducible nine-factor composite scoring
framework with SASB materiality weighting} for mid-cap ESG integration
(Section~3).
\item We provide \textbf{empirical evidence on ESG--financial integration in
mid-cap equities across the US and India}, using a common methodology across
both markets (Section~4).
\item We deliver \textbf{comprehensive robustness and overfitting control}
through bootstrap, subsampling, sector cross-validation, and walk-forward
evaluation, including strong stability diagnostics (CV ratio $=0.499$,
quasi-temporal ratio $=0.509$, walk-forward ratio $=0.978$) under
perturbation (Section~4).
\item We present a \textbf{transparent proxy-based ESG methodology with full
provenance tracking}, enabling auditability and replication with public data
(Section~3).
\end{itemize}

Empirically, the balanced top-20 portfolio generated by our framework delivers
a score-weighted cross-sectional information ratio of 0.475 (measuring score-return alignment rather than traditional realized-return IR), bear-market
excess of +7.27~percentage points with information ratio 0.585, and portfolio
beta of 0.82. In bear-market decomposition, excess momentum is +10.65~pp; in
rolling rebalancing, the strategy beats the benchmark (momentum proxy) in 100\% of periods with
average excess of +7.38\%. These performance metrics are based on
cross-sectional momentum proxies rather than realized portfolio returns and are
therefore evidence of ranking quality rather than realized investable alpha.
Factor diagnostics maintain low multicollinearity across factors (maximum VIF
= 2.80), PCA retains 3 components explaining 67.4\% of variance with
silhouette 0.261, bootstrap rank stability is perfect (Kendall $\tau=1.000$),
and generalization to large-cap universes is exceptionally strong: testing on
45 S\&P 500 companies yields Kendall $W=0.933$, 9/9 factors pass z-score
stability tests, and rank stability reaches Spearman $\rho=0.9993$. Monotonicity evidence is stronger
for the forward-quality proxy (7/8 factors) than for the return proxy, where
only 2/9 factors pass Jonckheere--Terpstra testing
(\texttt{risk\_adjusted}, $p<0.001$; \texttt{S\_score}, $p=0.036$). Capacity
analysis indicates the strategy is implementable from \$25M to \$500M, with
failure at \$1B+, transaction costs of 24--94~bps across AUM levels, and
25.4\% quarterly turnover. ESG appears primarily as a risk filter: ESG scores
are negatively associated with volatility ($\rho=-0.357$, $p<0.001$) and beta
($\rho=-0.146$, $p=0.015$). The observed ESG--financial link ($R^2=0.544$)
should not be read as an empirical discovery because it is partly induced by
the proxy-construction design. Factor ablation further identifies
\texttt{growth\_score} as the most influential component
($\tau=0.777$ when removed), and INCY (Healthcare) ranks first in the balanced
profile with preference score 85.4. We also explicitly distinguish between
cross-validated optimal weights and deployed weights: although CV optimization
selects a 77\% \texttt{market\_score} weight, the deployed balanced profile
uses 5\% to preserve interpretability and avoid over-concentration.
These results indicate that transparent ESG--financial integration can be both
practically competitive and methodologically testable in the mid-cap segment.
